\chapter{Mulch and evaporation layers}\label{ch:mulch}
\markboth{Mulch and evaporation}{Mulch and evaporation}
\textbf{After this chapter you can} compute what a mulch saves and where the saved water goes, follow the two stages of soil drying with the application's constants, and read the regime map.

\section{Mulch against bare soil}
The layer evaluates Engine 1 at bare soil ($\varphi=1$), at your mulch strength and at $\varphi=0.1$. With $E_\varphi,I_\varphi,D_\varphi$ from Equations 5.4, 5.7 and 5.9:
\begin{equation}\Delta E=E_{1.0}-E_\varphi,\quad\Delta I=I_{1.0}-I_\varphi,\quad\Delta D=D_\varphi-D_{1.0}.\end{equation}

\begin{example}{what a mulch saves}
For a soil the engine gives bare soil ($\varphi=1$): $E=190$, $I=180$, $D=40$; and with mulch $\varphi=0.5$: $E=140$, $I=148$, $D=62$ mm yr$^{-1}$ (illustration values). Find the three savings and the verdict; then repeat with a wet year ($\rho=1.3$) where $E$ falls from $175$ to $130$, $I$ from $120$ to $112$, and $D$ rises from $150$ to $188$.

\textbf{Solution.} Normal year: $\Delta E=190-140=50$ mm, $\Delta I=180-148=32$ mm, $\Delta D=62-40=22$ mm. The share of the saved evaporation that becomes less irrigation is $32/50=0.64$ and the share that becomes extra drainage is $22/50=0.44$ (the two exceed 1 together, so part of the drainage is not funded from the evaporation saving alone; the rest is storage change and transpiration, see Equation~(\ref{eq:fate2})). Since $\Delta I>\Delta D$ the verdict is ``dry to normal: most of the saved water becomes less irrigation''. Wet year: $\Delta E=45$, $\Delta I=8$, $\Delta D=38$. Now $\Delta D>\Delta I$ and the verdict reads ``wet conditions: much of the water the mulch saves now leaves as drainage''.
\end{example}

\section{Fate of the saved water}
At each rainfall scale $\rho$ the layer partitions the evaporation saved,
\begin{equation}\label{eq:fate2}f_I=\frac{\max\{\Delta I,0\}}{\max\{\Delta E,0.1\}},\ f_D=\frac{\max\{\Delta D,0\}}{\max\{\Delta E,0.1\}},\ f_T=\frac{\max\{\Delta T,0\}}{\max\{\Delta E,0.1\}},\ f_{\rm rest}=\max\{0,1-f_I-f_D-f_T\},\end{equation}
and reports as crossover the first $\rho$ with $\Delta D\ge\Delta I$ and $\Delta D>1$ mm.

\begin{example}{partition of the saved evaporation}
A mulch of strength $\varphi=0.5$ gives, at rainfall scales $\rho=0.5,\dots,1.4$, the changes of Table~\ref{tab:fatepart} (illustration values shaped like the study's trend: the irrigation saving falls with rain while drainage rises). Compute the partition of Equation~(\ref{eq:fate2}) and the crossover.

\textbf{Solution.} At $\rho=1.0$: $f_I=26/50=0.52$, $f_D=13/50=0.26$, $f_T=4/50=0.08$, $f_{\rm rest}=1-0.86=0.14$. The other rows are in the table.\begin{table}[H]\centering\small\caption{Partition of the saved evaporation.}\label{tab:fatepart}
\begin{tabular}{lrrrrrrrr}\toprule
$\rho$ & $\Delta E$ & $\Delta I$ & $\Delta D$ & $\Delta T$ & $f_I$ & $f_D$ & $f_T$ & $f_{\rm rest}$\\\midrule
0.5 & 38 & 33 & 0.6 & 1 & 0.87 & 0.02 & 0.03 & 0.09\\
0.7 & 44 & 31 & 3 & 2 & 0.70 & 0.07 & 0.05 & 0.18\\
0.9 & 48 & 29 & 8 & 3 & 0.60 & 0.17 & 0.06 & 0.17\\
1.0 & 50 & 26 & 13 & 4 & 0.52 & 0.26 & 0.08 & 0.14\\
1.1 & 50 & 21 & 17 & 4 & 0.42 & 0.34 & 0.08 & 0.16\\
1.2 & 49 & 15 & 24 & 4 & 0.31 & 0.49 & 0.08 & 0.12\\
1.3 & 48 & 9 & 29 & 5 & 0.19 & 0.60 & 0.10 & 0.10\\
1.4 & 46 & 5 & 33 & 5 & 0.11 & 0.72 & 0.11 & 0.07\\\bottomrule\end{tabular}\end{table}

The first $\rho$ in the table with $\Delta D\ge\Delta I$ and $\Delta D>1$ is $\rho=1.2$. Between the grid points $\Delta D-\Delta I$ is $-4$ at $\rho=1.1$ and $+9$ at $\rho=1.2$, so linear interpolation puts the crossover at $\rho=1.1+0.1\times4/13=1.131$. The study reports about 60 per cent of the saved water going to reduced irrigation in dry years, about 63 per cent to extra drainage in wet years, and a crossover near $\rho\approx1.2$; the illustration table, which is not the study's, gives $f_I=0.87$ at $\rho=0.5$ and $f_D=0.72$ at $\rho=1.4$, which have the same trend.
\end{example}

\section{The regime map}
For each cell of (mulch factor, rainfall scale) the engine forms the regime index $\mathcal I=(E-D)/\max\{1,E+D\}$: evaporation-dominated if $\mathcal I>0.05$, drainage-dominated if $\mathcal I<-0.05$, on the boundary otherwise; $\mathcal I=+1$ is all evaporation and $-1$ all drainage.

\begin{example}{regime index}
Classify the seven cells of Table~\ref{tab:regime} (illustration values of annual $E$ and $D$).

\textbf{Solution.} Compute $\mathcal I=(E-D)/\max\{1,E+D\}$. For the first cell: $(190-12)/202=+0.881$, evaporation-dominated. For the cell $(\varphi=0.5,\rho=1.0)$: $(140-62)/202=+0.386$. For $(0.5,1.3)$: $(130-188)/318=-0.182$, drainage-dominated. For $(0.2,1.0)$: $(70-95)/165=-0.152$.\begin{table}[H]\centering\small\caption{Regime index.}\label{tab:regime}
\begin{tabular}{lrrrrr}\toprule
$\varphi$ & $\rho$ & $E$ & $D$ & $\mathcal I$ & Regime\\\midrule
1.0 & 0.7 & 190 & 12 & +0.881 & evaporation\\
1.0 & 1.0 & 185 & 45 & +0.609 & evaporation\\
1.0 & 1.3 & 175 & 150 & +0.077 & evaporation\\
0.5 & 0.7 & 125 & 15 & +0.786 & evaporation\\
0.5 & 1.0 & 140 & 62 & +0.386 & evaporation\\
0.5 & 1.3 & 130 & 188 & -0.182 & drainage\\
0.2 & 1.0 & 70 & 95 & -0.152 & drainage\\\bottomrule\end{tabular}\end{table}
\textbf{Reading.} Strong mulch in a wet climate pushes the soil into the drainage regime where mulch ``saves less'' water for the crop, because the saved evaporation leaves as drainage.
\end{example}

\section{Two-stage evaporation}
With $e_0$ the atmospheric demand under mulch and canopy, $Q_1$ the stage-1 storage and $\tau$ the stage-2 time scale, $t_1=(Q_1-q_0)/e_0$ and
\begin{equation}Q(t)=\begin{cases}q_0+e_0t,&t\le t_1\\ Q_1+2e_0\tau\bigl(\sqrt{1+(t-t_1)/\tau}-1\bigr),&t>t_1\end{cases}\qquad e(t)=\begin{cases}e_0,&t\le t_1\\ e_0/\sqrt{1+(t-t_1)/\tau},&t>t_1\end{cases}\end{equation}
The parameters are
$h_i=\frac{0.25}{\alpha}(1.5-0.5\rho)$, $h_s=4/\alpha$, $Q_1=\max\{2,[\theta(h_i)-\theta(h_s)]L_1(0.6+0.4\rho)\}$ with $L_1=40$ mm, $\tau=0.006+0.03(\mathrm{clay}+\mathrm{silt})/100$ d, $q_0=(1-w_f)0.95Q_1$, $w_f=\mathrm{clip}(1-z_d/50;0.05,1)$, and $e_0=e_{0,b}\varphi(1-0.4c)$.

\begin{example}{stage-1 and stage-2 evaporation of the ten soils}
With bare soil ($\varphi=1$, $c=0$, so $e_0=4.2$ mm d$^{-1}$), rainfall scale $\rho=1$ and a surface method ($q_0=0$), compute for the ten soils $h_i$, $h_s$, $Q_1$, the stage-1 duration $t_1$, $\tau$ and the cumulative loss at 10 days.

\textbf{Solution, loam S06 in full.} $\alpha=0.036$: $h_i=0.25/0.036\times(1.5-0.5)=6.94$ cm; $h_s=4/0.036=111$ cm. The retention uses $\phi=0.428$, $\theta_r=0.022$ (Chapter~\ref{ch:struct}): $\Se(6.94)=0.9959$$\Se(6.94)=0.9617$ and $\Se(111)=0.4427$, so $\theta(h_i)-\theta(h_s)=0.406(0.9617-0.4427)=0.2107$ and $Q_1=8.43$ mm. Stage 1 lasts $t_1=Q_1/e_0=2.01$ d. Then $\tau=0.006+0.03(0.60)=0.024$ d. After 10 days: $t-t_1=7.99$ d, $Q=Q_1+2(4.2)(0.024)(\sqrt{1+333.0}-1)=11.9$ mm. Table~\ref{tab:stage} gives all ten soils.\begin{table}[H]\centering\small\caption{Stage-1 storage, duration and cumulative loss in 10 days (bare, $\rho=1$).}\label{tab:stage}
\begin{tabular}{lrrrrrrr}\toprule
ID & Soil & $h_i$ & $h_s$ & $Q_1$ (mm) & $t_1$ (d) & $\tau$ (d) & $Q(10)$ (mm)\\\midrule
S01 & Coarse loam & 2.6 & 42 & 5.99 & 1.43 & 0.0225 & 9.5\\
S02 & Coarse sand & 4.7 & 75 & 7.85 & 1.87 & 0.0096 & 10.1\\
S03 & Sand & 1.7 & 28 & 12.71 & 3.03 & 0.0084 & 14.7\\
S04 & Loamy sand & 2.0 & 32 & 11.90 & 2.83 & 0.0114 & 14.2\\
S05 & Sandy loam & 3.3 & 53 & 10.49 & 2.50 & 0.0165 & 13.3\\
S06 & Loam & 6.9 & 111 & 8.43 & 2.01 & 0.0240 & 11.9\\
S07 & Silt loam & 12.5 & 200 & 7.07 & 1.68 & 0.0300 & 11.0\\
S08 & Sandy clay loam & 4.2 & 68 & 7.35 & 1.75 & 0.0195 & 10.6\\
S09 & Clay loam & 13.2 & 211 & 5.62 & 1.34 & 0.0270 & 9.5\\
S10 & Silty clay loam & 25.0 & 400 & 4.60 & 1.10 & 0.0330 & 8.9\\\bottomrule\end{tabular}\end{table}
\textbf{Check against the manual.} It says stage 1 lasts about 3 days for the sand at the season-mean demand and about 1 day for the silty clay loam: the table gives $3.03$ and $1.10$ days. Stage 1 is longest in the sands (large $Q_1$, because the window between $h_i$ and $h_s$ scales with $1/\alpha$) while in stage 2 the rate $e_0/\sqrt{1+(t-t_1)/\tau}$ falls faster for the small-$\tau$ sands than for the fine soils: the sands lose their larger stage-1 reserve first and then dry out quickly, as the application describes.
\end{example}

\begin{example}{mulch and canopy in the evaporation model}
For the loam, with the stage constants above, compute $e_0$, $t_1$, and the cumulative loss at 10 and 24 days for mulch strengths $\varphi=1,0.5,0.2$ with canopy cover $c=0$ and $c=0.8$.

\textbf{Solution.} $e_0=4.2\varphi(1-0.4c)$. For $\varphi=0.5$, $c=0.8$: $e_0=4.2\times0.5\times0.68=1.428$ mm d$^{-1}$, $t_1=8.43/1.428=5.90$ d. Table~\ref{tab:mulchev} lists the six cases. Mulch lengthens stage 1 in proportion to $1/\varphi$ (from 2.0 days to 10 days for $\varphi=0.2$), and the cumulative loss falls more than in proportion because stage 2 starts later, which is why the evaporation saving grows convexly with mulch strength.\begin{table}[H]\centering\small\caption{Effect of mulch and canopy on evaporation of the loam.}\label{tab:mulchev}
\begin{tabular}{lrrrrr}\toprule
$\varphi$ & $c$ & $e_0$ & $t_1$ (d) & $Q(10)$ & $Q(24)$\\\midrule
1.0 & 0 & 4.200 & 2.01 & 11.9 & 14.3\\
1.0 & 0.8 & 2.856 & 2.95 & 10.6 & 12.4\\
0.5 & 0 & 2.100 & 4.01 & 9.9 & 11.2\\
0.5 & 0.8 & 1.428 & 5.90 & 9.3 & 10.2\\
0.2 & 0 & 0.840 & 10.03 & 8.4 & 9.4\\
0.2 & 0.8 & 0.571 & 14.76 & 5.7 & 8.9\\\bottomrule\end{tabular}\end{table}
\end{example}

\begin{example}{the drying front}
After stage 1, the depth of the dry layer is $\ell(t)=\min\{60,(Q(t)-Q_1)/\Delta\theta\}$ mm with $\Delta\theta=\mathrm{clip}(\theta(h_i)-\theta(10^4);0.05,0.4)$. For the loam bare soil at 10 days, find $\ell$ and the suction at depth $z=40$ mm from $\log_{10}h(z)=\log_{10}h_s+(\log_{10}h_i-\log_{10}h_s)(1-e^{-(z-\ell)/\lambda})$, $\lambda=L_1/3$.

\textbf{Solution.} $\Delta\theta=\theta(6.94)-\theta(10^4)=0.375$. At 10 days $Q=11.9$ mm, so $\ell=(11.9-8.43)/0.375=9.3$ mm. With $\lambda=13.3$ mm, at $z=40$ mm the exponent is $-(40-9.3)/13.3=-2.30$ and $\log_{10}h=0.962$, that is $h=9$ cm.
\end{example}
